\section{Purpose}

\textbf{[STIPULATION]} In FFGFT all formulas are set up in natural units
($c=\hbar=k_B=1$). When back-computing to SI quantities, the physical constants
reappear with specific powers. These powers are \emph{not} freely choosable --
they follow necessarily from dimensional analysis. An error in the power of $c$
is an error of factors $\sim3\times10^8$ to $(3\times10^8)^4$ -- not noticeable
when looking only at orders of magnitude, fatal when comparing with a measurement.

This document gives the systematic procedure: first dimensional analysis in
natural units, then automated back-computation, then verification. It is not a
teaching document on unit systems (that is A080/A085) -- it is a practical
toolbox for concrete calculations.

\section{Basic Principle: All Dimensions as Energy Powers}

\textbf{[B]} In natural units $c=\hbar=k_B=1$. Every physical quantity thereby
has a unique dimension $[E]^n$ with rational $n$:

\begin{center}
\renewcommand{\arraystretch}{1.35}
{\small
\begin{tabular}{llll}
\toprule
Quantity & Nat.\ dim. & SI origin & $n$ \\
\midrule
Energy & $[E]^1$ & --- & $+1$ \\
Mass & $[E]^1$ & $E=mc^2\Rightarrow m=E/c^2$ & $+1$ \\
Momentum & $[E]^1$ & $p=E/c$ & $+1$ \\
Length & $[E]^{-1}$ & $\ell=\hbar c/E$ & $-1$ \\
Time & $[E]^{-1}$ & $t=\hbar/E$ & $-1$ \\
Velocity & $[E]^0$ & $v=\ell/t=c=1$ & $0$ \\
Force & $[E]^2$ & $F=E/\ell=E^2/(\hbar c)$ & $+2$ \\
Pressure/energy density & $[E]^4$ & $\rho=E/\ell^3=E^4/(\hbar c)^3$ & $+4$ \\
Charge & $[E]^0$ & from $\alpha=e^2/(4\pi\varepsilon_0\hbar c)$ dimensionless & $0$ \\
Temperature & $[E]^1$ & $T=E/k_B$ & $+1$ \\
Entropy & $[E]^0$ & $S=k_B\cdot(\text{dimensionless})$ & $0$ \\
Gravitational constant $G$ & $[E]^{-2}$ & $G=\ell_P^2c^3/\hbar$ & $-2$ \\
\bottomrule
\end{tabular}
}
\renewcommand{\arraystretch}{1.0}
\end{center}

\textbf{[B]} \textbf{Rule.} From $n$ the back-computation factor follows
directly: a quantity with $[E]^n$ in natural units has the SI value
\[
  X_\text{SI} = X_\text{nat}\cdot\left(\frac{\hbar c}{1\,\text{SI-length}}\right)^n.
\]
In practice: assemble the SI factor from $\hbar$, $c$, $k_B$ so that the SI
dimension comes out.

\section{The $c$-Power in Back-Computation: Systematic Derivation}

\textbf{[B]} The question ``which power of $c$?'' is answered from the SI
dimension of the target quantity. The method in three steps:

\textbf{Step 1: Write down the SI target dimension.}
E.g.\ energy density: $[\rho]=\text{J/m}^3=\text{kg\,m}^{-1}\text{s}^{-2}$.

\textbf{Step 2: Express SI dimension through $\hbar^a c^b k_B^d$.}
$\hbar$ has SI dim $[\text{J\,s}]$, $c$ has $[\text{m/s}]$, $k_B$ has $[\text{J/K}]$.
Find $a,b,d\in\mathbb Z$ such that $[\hbar^a c^b k_B^d]=$ target dim$/[E_\text{nat}^n]$.

\textbf{Step 3: Verification by substitution.}
Insert the found $(a,b,d)$ and check against the measured SI quantity.

\textbf{[K]} \textbf{Table of most common back-computation factors:}

\begin{center}
\renewcommand{\arraystretch}{1.35}
{\small\setlength\tabcolsep{4pt}
\begin{tabular}{lllp{4.5cm}}
\toprule
Quantity & Nat.\ unit & SI factor & Derivation \\
\midrule
Mass & $[E]$ in MeV & $\div c^2$ & $m_e=0{.}511\,\text{MeV}/c^2$ \\
Length & $[E]^{-1}$ & $\times\hbar c$ & $\ell_P=1/E_P\cdot\hbar c$ \\
Time & $[E]^{-1}$ & $\times\hbar$ & $t_P=1/E_P\cdot\hbar$ \\
Velocity & dimensionless & $\times c$ & $v=0{.}5\cdot c$ \\
Force & $[E]^2$ & $\div(\hbar c)$ & $F[\text{N}]=F_\text{nat}\cdot E^2/(\hbar c)$ \\
Energy density & $[E]^4$ & $\div(\hbar c)^3$ & $\rho[\text{J/m}^3]=\rho_\text{nat}\cdot E^4/(\hbar c)^3$ \\
Temperature & $[E]$ & $\div k_B$ & $T[\text{K}]=E[\text{J}]/k_B$ \\
Frequency $f$ & $[E]^1$ & $\div\hbar$ & $f[\text{s}^{-1}]=E[\text{J}]/\hbar$ \\
$G$ & $[E]^{-2}$ & $\times\hbar c^3$ & $G[\text{m}^3/\text{kg\,s}^2]=G_\text{nat}\cdot\hbar c^3$ \\
\bottomrule
\end{tabular}
}
\renewcommand{\arraystretch}{1.0}
\end{center}

\section{Concrete Examples from FFGFT}

\subsection{Mass: One Power of $c^2$}

\textbf{[K]} From $E=mc^2$ it follows that $m=E/c^2$. A value in natural units
(MeV) becomes a SI mass by dividing by $c^2$:
\[
  m_e = 0{.}511\,\frac{\text{MeV}}{c^2}
  = \frac{0{.}511\times10^6\times1{.}602\times10^{-19}\,\text{J}}{(2{.}998\times10^8\,\text{m/s})^2}
  = 9{.}109\times10^{-31}\,\text{kg}.
\]
Error $c^1$ instead of $c^2$: factor $3\times10^8$ off.
Error $c^0$ (no $c$): factor $(3\times10^8)^2\approx10^{17}$ off.

\subsection{Planck Length: $\hbar c$ Together, Not $c$ Alone}

\textbf{[K]} $\ell_P$ in natural units is an inverse energy. Back-computation to
metres needs $\hbar c$:
\[
  \ell_P = \frac{\hbar c}{E_P}
  = \frac{(1{.}055\times10^{-34}\,\text{J\,s})(2{.}998\times10^8\,\text{m/s})}
         {1{.}956\times10^9\,\text{J}}
  = 1{.}616\times10^{-35}\,\text{m}.
\]
Only $\hbar/E_P$ (without $c$) gives the Planck \emph{time}, not the Planck
\emph{length} -- a typical error.

\subsection{Energy to Inverse Time: Only $\hbar$, No $c$}

\textbf{[K]} An energy $E$ (dimension $[E]^1$) is back-computed to inverse time
$[\text{s}^{-1}]$ by dividing by $\hbar$:
\[
  f = \frac{E}{\hbar},
  \qquad
  [f] = \frac{[\text{J}]}{[\text{J\,s}]} = [\text{s}^{-1}].\checkmark
\]
Example: $E=1{.}41\times10^{-33}\,\text{eV}=2{.}26\times10^{-52}\,\text{J}$ gives
$f=E/\hbar=2{.}14\times10^{-18}\,\text{s}^{-1}$ (A265). Dividing by $\hbar c$
instead gives an inverse length (wrong dimension). The $c$-power is zero -- it
does not appear.

\subsection{Energy Density: Four Powers of $c$ in the Denominator}

\textbf{[K]} Casimir energy density $\rho$ has dimension $[E]^4$ in natural
units. Back-computation:
\[
  \rho[\text{J/m}^3]
  = \rho_\text{nat}\,[E^4]\cdot\frac{1}{(\hbar c)^3}.
\]
Three powers of $\hbar c$ in the denominator -- because length contributes
$[E]^{-1}$ three times. $c$ appears as $c^3$ (from the three $\hbar c$).

\subsection{Gravitational Constant: $c^3$ in the Numerator}

\textbf{[K]} In natural units $G$ has dimension $[E]^{-2}$. The SI constant
$G=6{.}674\times10^{-11}\,\text{m}^3\text{kg}^{-1}\text{s}^{-2}$ requires:
\[
  G_\text{SI} = G_\text{nat}\cdot\hbar c^3,
  \qquad
  [G_\text{SI}] = [E]^{-2}\cdot[\hbar][c]^3
  = \frac{\text{J\,s}\cdot\text{m}^3/\text{s}^3}{\text{J}^2}
  = \frac{\text{m}^3}{\text{kg\,s}^2}.\checkmark
\]
Error: $c^2$ instead of $c^3$ gives wrong dimension immediately recognisable.

\section{The General Procedure: Dimension Matrix}

\textbf{[B]} For every formula $X=f(\xi,E_0,\ldots)$ in natural units:

1. Determine $n$ from the table above (or by counting lengths/masses/times in
   the SI dimension of $X$).

2. Write the SI factor:
\[
  X_\text{SI} = X_\text{nat}\cdot\hbar^a\cdot c^b\cdot k_B^d
\]
with $a,b,d$ from the linear system:
\begin{align*}
  [\text{kg}]:&\quad -a+b = n_\text{kg} \\
  [\text{m}]:&\quad a+b = n_\text{m} \\
  [\text{s}]:&\quad -a-b = n_\text{s}
\end{align*}
(where $n_\text{kg}$, $n_\text{m}$, $n_\text{s}$ are the exponents of the SI
dimension of $X$; $k_B$ is added when kelvin appears.)

3. Verification: insert $\hbar=1{.}055\times10^{-34}\,\text{J\,s}$,
   $c=2{.}998\times10^8\,\text{m/s}$,
   $k_B=1{.}381\times10^{-23}\,\text{J/K}$ and check the SI unit.

\textbf{[K]} \textbf{When $k_B$ appears and when not.} $k_B$ appears exactly when
the target quantity contains kelvin (temperature, entropy/kelvin) or when
converting from an energy to a kelvin quantity. In all other cases (masses,
lengths, times, forces) $k_B$ does not appear.

\section{Unit Check of a Formula}

\textbf{[B]} Every formula set up in FFGFT should pass the unit check: both sides
of the equation must have the same energy power $[E]^n$ in natural units, and
on back-computation the same SI unit must emerge.

\textbf{[K]} \textbf{Example: $T\cdot E = 1$} (A060). In natural units:
$[T]=[E]^{-1}$ (time) and $[E]=[E]^{+1}$. Product: $[E]^0$ -- dimensionless.
The product $T\cdot E=1$ is a correct equation between a time and an inverse time
-- the unit check is passed.

\textbf{[K]} \textbf{Example: $\alpha=\xi\cdot E_0^2$ -- unit check of the
$\alpha$ formula.} The formula is \emph{correct}. The unit check shows why --
and simultaneously makes the implicit reference scale visible.

$\alpha$ is dimensionless, $\xi$ is dimensionless. The equation therefore holds
only if $E_0^2$ is also dimensionless. That is the case when $E_0$ appears as a
pure numerical ratio -- i.e.\ measured in MeV and divided by $1\,\text{MeV}$:
\[
  \alpha = \xi\cdot\Bigl(\frac{E_0}{1\,\text{MeV}}\Bigr)^2.
\]
The $1\,\text{MeV}$ in the denominator is the \emph{implicit reference scale} --
it makes $E_0/1\,\text{MeV}$ dimensionless and the equation correct. It does not
stand explicitly in the short notation, but is indispensable.

\textbf{Numerical verification:}
\[
  \alpha = \frac{4}{30000}\cdot\Bigl(\frac{7{.}398\,\text{MeV}}{1\,\text{MeV}}\Bigr)^2
  = 1{.}333\times10^{-4}\times54{.}73
  = 7{.}297\times10^{-3}
  = \frac{1}{137{.}0}.
\]
Agreement with experiment to $0{.}0005\,\%$ -- the formula is correct.

\textbf{What the unit check reveals:} the choice of MeV as reference scale is not
free -- it is the empirical SI anchor of the FFGFT sector. A different reference
scale (e.g.\ GeV) would give $\xi$ a different numerical value, but $\alpha$
remains the same physical number. This is the statement from A261: $E_0=7{.}398\,\text{MeV}$
carries the anchor, and therefore one may not simply insert different energy units
in this formula without simultaneously changing $\xi$.

\textbf{[S]} \textbf{Converting to natural units.} Translating a SI formula to
natural units: replace all $\hbar$, $c$, $k_B$ by $1$ and drop the units. The
resulting formula is then valid only in natural units; on back-computation all
factors return.

\section{Symbol Glossary}

Symbols used throughout the entire A-series are explained in A010.
Specific to this document:

\begin{center}
\renewcommand{\arraystretch}{1.3}
{\small\begin{tabular}{lp{9cm}}
\toprule
Symbol & Meaning \\
\midrule
$[X]$ & Dimension of quantity $X$ \\
$[E]^n$ & Energy dimension to power $n$ \\
$n_\text{kg},n_\text{m},n_\text{s}$ & Exponents of the SI unit of $X$
  in $\text{kg}^{n_\text{kg}}\text{m}^{n_\text{m}}\text{s}^{n_\text{s}}$ \\
$a,b,d$ & Exponents of $\hbar^a c^b k_B^d$ in the back-computation factor \\
$E_P = \sqrt{\hbar c^5/G}\approx1{.}956\times10^9\,\text{J}$ & Planck energy \\
\bottomrule
\end{tabular}}
\renewcommand{\arraystretch}{1.0}
\end{center}

\section{Verification Script}

\begin{itemize}[leftmargin=2em,itemsep=2pt,topsep=2pt]
  \item Numerical verification of all back-computation examples; dimension matrix
    solver for arbitrary SI target dimensions; verification for $m_e$, $\ell_P$,
    inverse time, $\rho_\text{Casimir}$, $G$:
    \texttt{python/A\_Serie\_Skripte/a266\_einheitenpruefung.py}
\end{itemize}

\section{Open Questions}

\textbf{The dimension matrix presupposes that the formula is set up in pure Planck
units.} Mixed systems (e.g.\ MeV for energy, fm for length) require additional
conversion steps.

\textbf{Non-integer energy powers.} Some FFGFT expressions have rational $n$
(e.g.\ charge $[E]^{1/2}$). The procedure applies in principle, but the linear
system then gives rational solutions.

\section{Supersedes}

This document subsumes: Dok.~013 (dimensional analysis, gravitational constant
from $\xi$, Planck length, conversion factors), Dok.~014 (natural units,
dimension table, scaling factor $S_{T0}$, derivation vs.\ back-computation),
Dok.~015 (unit conversions, energy dimension table, force/pressure/charge).
\emph{Not} taken over: pedagogical derivation of the individual constants from
$\xi$ (A015, A145) and the unit convention as such (A080, A085).

\section{Use of AI Tools}

This document was created with the assistance of AI tools. Responsibility
for content and errors rests with the author. Full wording of the rules in A010.
